\Askhsh{23314_SOLUTION}{1}
\begin{ylh}{1}
    \item [Ύλη από εικόνα]
\end{ylh}
ΛΥΣΗ
\begin{alist}
    \item[α)]
    \begin{rlist}
        \item[i)] $\lim_{x \to 0} f(x) = 2$
        \item[ii)] $\lim_{x \to -2^+} f(x) = 0$
        \item[iii)] $\lim_{x \to -2^-} f(x) = 0$
    \end{rlist}

    \item[β)]
    \begin{rlist}
        \item[i)] Επειδή $\lim_{x \to -2^+} f(x) = 0$, με $f(x) < 0$, καθώς το $x$ προσεγγίζει το $-2$ από μεγαλύτερες τιμές ($x>-2$), έχουμε $\lim_{x \to -2^+} \frac{1}{f(x)} = -\infty$.
        \item[ii)] Αν θέσουμε $u = f(x)$, τότε $\lim_{x \to -2^-} u = \lim_{x \to -2^-} f(x) = 0$, με $f(x) > 0$ καθώς το $x$ προσεγγίζει το $-2$ από μικρότερες τιμές ($x<-2$) και $\lim_{u \to 0^+} \ln u = -\infty$, επομένως $\lim_{x \to -2^-} \ln(f(x)) = -\infty$.
    \end{rlist}
\end{alist}